\newpage
\subsection{Общая структурная схема электронных часов}
\label{sec:clock-structure}

Общая структурная схема приведена на рис.~\ref{fig:clock-structure}. Устройство управления формирует импульсные команды $i_j$ и потенциальные сигналы $L_k$. Шины $T$, $B$ и $A$ передают данные текущего времени, секундомера и будильника соответственно. Точка на линии обозначает соединение.

\begin{figure}[H]
	\centering
	\begin{tikzpicture}[
		block/.style={draw,align=center,font=\small,minimum height=11mm,inner sep=3pt},
		wire/.style={-{Stealth[length=2mm]}},
		bus/.style={wire,semithick},
		signal/.style={font=\footnotesize,inner sep=2pt,fill=white},
	]
		\node[block,minimum width=2.4cm] (keys) at (1.9,0)
			{Кнопки $a,b,c$};
		\node[block,minimum width=4cm] (clock) at (12.5,0)
			{Генератор и делитель};

		\node[block,minimum width=3.4cm] (encoder) at (1.9,-2.05)
			{Преобразователь\\входных событий};
		\node[block,minimum width=3.7cm] (controller) at (7.05,-2.05)
			{Устройство\\управления};
		\node[block,minimum width=3.6cm] (timer) at (13.3,-2.05)
			{Таймер\\сигнализации};

		\node[block,minimum width=3.35cm] (time) at (1.9,-5.65)
			{Блок текущего\\времени};
		\node[block,minimum width=3.65cm] (stopwatch) at (6.85,-5.65)
			{Секундомер\\с триггером $L_9$};
		\node[block,minimum width=4.25cm] (alarm) at (12.65,-5.65)
			{Будильник\\с триггером $L_{10}$};

		\node[block,minimum width=5.1cm] (display) at (5.45,-8.55)
			{Мультиплексор и\\индикаторный преобразователь};
		\node[block,minimum width=4cm] (sound) at (12.65,-8.55)
			{Звуковой блок};
		\node[block,minimum width=3.05cm] (indicator) at (5.45,-10.35)
			{Индикатор};
		\node[block,minimum width=3.05cm] (speaker) at (12.65,-10.35)
			{Излучатель звука};

		\draw[wire] (keys) -- (encoder);
		\draw[bus] (encoder) -- node[signal,above] {$x,s$} (controller);
		\draw[wire] (clock.west) -- node[signal,above] {$f_{\mathrm{clk}}$}
			(7.05,0) -- (controller.north);
		\draw[wire] (7.05,-0.77) -- (3.27,-0.77) -- (3.27,-1.5);
		\fill (7.05,-0.77) circle (1.3pt);
		\draw[wire] (13.3,-0.55) -- node[signal,right] {$f_{\mathrm{clk}}$} (timer.north);
		\draw[wire] (timer.east) -- (15.65,-2.05) -- (15.65,1.03) --
			node[signal,above] {$\rho$} (-0.34,1.03) -- (-0.34,-1.86) --
			(0.2,-1.86);
		\draw[wire] (time.west) -- (-0.34,-5.65) --
			node[signal,left] {$\tau$} (-0.34,-2.24) -- (0.2,-2.24);
		\draw[wire] (controller.east) -- node[signal,above] {$L_8$} (timer.west);
		\draw[wire] (9.6,-2.05) -- (9.6,-7.65) --
			node[signal,above] {$L_8$} (11.3,-7.65) -- (11.3,-8);
		\fill (9.6,-2.05) circle (1.3pt);

		\draw[bus] (5.6,-2.6) -- (5.6,-2.98) --
			node[signal,above,pos=0.31] {$i_1,i_2,i_3,L_1$} (2.7,-2.98) -- (2.7,-5.1);
		\draw[bus] (7.35,-2.6) -- node[signal,right,pos=0.53] {$i_4,i_5,i_6$}
			(7.35,-5.1);
		\draw[bus] (8.55,-2.6) -- (8.55,-3.15) --
			node[signal,above,pos=0.75] {$i_7,\ldots,i_{10}$} (12.65,-3.15) -- (12.65,-5.1);
		\draw[bus] (6.25,-2.6) -- (6.25,-3.62) -- (4.19,-3.62) --
			node[signal,left,pos=0.27] {$L_2,\ldots,L_7$} (4.19,-8);

		\draw[wire] (clock.east) -- (15.26,0) -- (15.26,-4.24) --
			node[signal,above,pos=0.14] {$1$ Гц} (1.12,-4.24) -- (1.12,-5.1);
		\draw[wire] (6.12,-4.24) -- (6.12,-5.1);
		\fill (6.12,-4.24) circle (1.3pt);

		\draw[bus] (1.9,-6.2) -- (1.9,-7.22) --
			node[signal,above] {$T$} (3.48,-7.22) -- (3.48,-8);
		\draw[bus] (6.2,-6.2) -- node[signal,right,pos=0.76] {$B$} (6.2,-8);
		\draw[bus] (11.1,-6.2) -- (11.1,-7.06) --
			node[signal,above] {$A$} (7.55,-7.06) -- (7.55,-8);
		\draw[bus] (3.05,-6.2) -- (3.05,-6.58) --
			node[signal,below,pos=0.22] {$T,p_{\mathrm M}$} (13.55,-6.58) -- (13.55,-6.2);
		\draw[bus] (14.1,-6.2) -- node[signal,right] {$p_{\mathrm A},n$} (14.1,-8);

		\draw[wire] (display) -- (indicator);
		\draw[wire] (sound) -- (speaker);
	\end{tikzpicture}
	\caption{Общая структурная схема электронных часов}
	\label{fig:clock-structure}
\end{figure}


Генератор задает опорную частоту $f_{\mathrm{clk}}$ для преобразователя входных событий, устройства управления и таймера сигнализации. Делитель формирует импульсы частотой 1~Гц для блока текущего времени и секундомера. Разрешения счета подаются раздельно, поэтому при останове часов в $S_{10}$ секундомер может продолжать отсчет.

\newpage
\subsubsection{Функции блоков}
\label{sec:clock-blocks}

Преобразователь входных событий устраняет дребезг контактов и регистрирует одно событие на нажатие. Он также принимает $\tau$ от блока текущего времени и $\rho$ от таймера. Каждому принятому событию соответствуют код $x$ (таблица~\ref{tab:inputs}) и синхроимпульс $s$ длительностью один такт. При отсутствии события выдается код $110$ при $s=0$.

Блок текущего времени содержит счетчики секунд и минут по модулю $60$ и счетчик часов по модулю $24$. При отсчете перенос передается из секунд в минуты и из минут в часы. Корректировка меняет только выбранное поле. При естественном переходе к новой минуте выдается импульс $p_{\mathrm M}$, а на минутах $00$, $15$, $30$ и $45$ также формируется событие $\tau$.

Блок секундомера содержит счетчики минут и секунд и триггер хода $L_9$. При $L_9=1$ секундные импульсы изменяют показание $B$ независимо от выбранного режима индикации.

Блок будильника хранит установленные часы и минуты, номер мелодии $n$ от 1 до 4 и признак включения $L_{10}$. По импульсу $p_{\mathrm M}$ блок сравнивает установленное и текущее время с точностью до минуты. При совпадении и $L_{10}=1$ импульс $p_{\mathrm A}$ поступает в звуковой блок и запускает воспроизведение мелодии $n$.

Мультиплексор выбирает источник данных по сигналам устройства управления. Индикаторный преобразователь выводит цифры в разрешенных полях. Признак включения будильника поступает по шине $A$.

При входе в $S_{12}$ сигнал $L_8$ включает четвертьчасовой звук и запускает таймер. Таймер отсчитывает одну секунду по опорным тактам и выдает событие $\rho$. Автомат возвращается в $S_0$, после чего $L_8$ снимается. Звуковой блок подает четвертьчасовой сигнал и мелодию будильника на общий излучатель.

\newpage
\subsubsection{Устройство управления}
\label{sec:clock-unit}

Согласно структуре из методических указаний (рис.~16 и 22) устройство управления включает функциональный преобразователь $F$, блок элементов памяти ЭП, формирователь $F_i$ и дешифратор $F_L$ (рис.~\ref{fig:clock-controller}).

\begin{figure}[H]
	\centering
	\begin{tikzpicture}[
		block/.style={draw,align=center,font=\small,minimum width=2.8cm,minimum height=10mm,inner sep=3pt},
		wire/.style={-{Stealth[length=2mm]}},
		signal/.style={font=\footnotesize,inner sep=2pt,fill=white},
	]
		\node[block] (f) at (3.75,0) {Преобразователь\\$F$};
		\node[block] (memory) at (9.35,0) {Память состояния\\ЭП};
		\node[block] (pulse) at (3.75,-2.65) {Формирователь\\$F_i$};
		\node[block] (level) at (9.35,-2.65) {Дешифратор\\$F_L$};

		\draw[wire] (0.42,0) -- node[signal,above] {$x$} (f);
		\draw[wire] (f) -- node[signal,above] {$q^+$} (memory);
		\draw[wire] (12.8,0) -- node[signal,above] {$f_{\mathrm{clk}}$} (memory.east);
		\draw[wire] (memory.north) -- (9.35,1.15) --
			node[signal,above] {$q$} (3.75,1.15) -- (f.north);
		\draw[wire] (f) -- node[signal,right,pos=0.28] {$y$} (pulse);
		\draw[wire] (9.35,-0.54) -- node[signal,right] {$q$} (level);
		\draw[wire] (0.42,-1.41) -- node[signal,above] {$s$}
			(2.05,-1.41) -- (2.05,-2.65) -- (pulse);
		\draw[wire] (2.05,-1.41) -- (7.14,-1.41) -- (7.14,-0.28) -- (memory.west |- 0,-0.28);
		\fill (2.05,-1.41) circle (1.3pt);
		\draw[wire] (pulse.south) -- node[signal,right] {$i_1,\ldots,i_{10}$} ++(0,-0.78);
		\draw[wire] (level.south) -- node[signal,right] {$L_1,\ldots,L_8$} ++(0,-0.78);
	\end{tikzpicture}
	\caption{Структура устройства управления}
	\label{fig:clock-controller}
\end{figure}


Блок ЭП хранит четырехразрядный код текущего состояния $q$ из таблицы~\ref{tab:states}. Преобразователь $F$ по кодам $q$ и $x$ вычисляет код следующего состояния $q^+$ и четырехразрядный код выхода $y$ по таблице~\ref{tab:binary-functions} и правилам для событий времени. Дешифратор $F_L$ непрерывно формирует сигналы $L_1,\ldots,L_8$ по коду $q$.

При $s=1$ на одном фронте опорного такта ЭП записывает $q^+$, а исполнительные блоки принимают импульсные команды. Таким образом, команда определяется состоянием, из которого совершается переход. При включении питания в ЭП записывается код $S_0$, а признаки хода секундомера и включения будильника сбрасываются.

\newpage
\subsubsection{Импульсные микрокоманды}
\label{sec:clock-pulses}

Импульсная микрокоманда вызывает одно действие при обработке события. Она может выдаваться и на петле графа, когда состояние сохраняется. Формирователь $F_i$ декодирует выход $z_j$ и разрешает команду только на время синхроимпульса $s$:
\[
	i_j = s\,[y=\operatorname{code}(z_j)], \qquad j=1,\ldots,10,
\]
где $[P]$ равен единице при истинном условии $P$ и нулю в остальных случаях, а $\operatorname{code}(z_j)$~--- двоичный код выхода из таблицы~\ref{tab:outputs}. Для $z_0$ все $i_j$ равны нулю. Назначения команд приведены в таблице~\ref{tab:clock-pulses}.

\begin{table}[H]
	\centering
	\caption{Импульсные микрокоманды устройства управления}
	\label{tab:clock-pulses}
	\small\setstretch{1.12}
	\renewcommand{\arraystretch}{1.12}
	\begin{tabularx}{\textwidth}{|c|>{\raggedright\arraybackslash}p{3.25cm}|>{\raggedright\arraybackslash}X|}
		\hline
		Команда & Блок & Действие \\
		\hline
		$i_1$ & Текущее время & Сбросить секунды до $00$. \\
		$i_2$ & Текущее время & Увеличить минуты на единицу по модулю $60$. \\
		$i_3$ & Текущее время & Увеличить часы на единицу по модулю $24$. \\
		$i_4$ & Секундомер & Установить $L_9$ и начать или продолжить отсчет. \\
		$i_5$ & Секундомер & Сбросить $L_9$ и зафиксировать достигнутое показание. \\
		$i_6$ & Секундомер & Обнулить показание остановленного секундомера. \\
		$i_7$ & Будильник & Инвертировать сохраняемый признак включения $L_{10}$. \\
		$i_8$ & Будильник & Увеличить заданные минуты на единицу по модулю $60$. \\
		$i_9$ & Будильник & Увеличить заданные часы на единицу по модулю $24$. \\
		$i_{10}$ & Будильник & Перейти к следующей мелодии. После четвертой выбирается первая. \\
		\hline
	\end{tabularx}
\end{table}

При нажатии $b$ в $S_6$ выходу $z_6$ соответствует импульс $i_6$, обнуляющий показание секундомера. При нажатии $c$ в $S_5$ выдается $z_0$, и триггер хода сохраняет значение. Дешифратор $F_L$ переключает индикацию при переходе в $S_0$.

\newpage
\subsubsection{Потенциальные микрокоманды}
\label{sec:clock-levels}

Потенциальные сигналы дешифратора $F_L$ сохраняют значения в течение выбранного режима.

\begin{itemize}
	\item $L_1$ разрешает ход счетчика текущего времени. В $S_{10}$ этот сигнал равен нулю.
	\item Пара $L_2L_3$ выбирает данные для индикации. Код $00$ соответствует текущему времени $T$, $01$~--- показанию секундомера $B$, $10$~--- времени будильника, $11$~--- номеру мелодии из шины $A$.
	\item $L_4$, $L_5$ и $L_6$ включают левое, среднее и правое поля. Секундомер использует среднее и правое, номер мелодии~--- левое.
	\item $L_7$ разрешает индикацию включенного будильника при $L_{10}=1$.
	\item $L_8$ разрешает звук в $S_{12}$. Его фронт запускает таймер.
\end{itemize}

В таблице~\ref{tab:clock-levels} режимы $S_4$, $S_5$ и $S_6$ имеют одинаковые сигналы индикации. Ход секундомера определяется уровнем $L_9$.

\begin{table}[H]
	\centering
	\caption{Потенциальные микрокоманды дешифратора $F_L$}
	\label{tab:clock-levels}
	\small\setstretch{1.12}
	\renewcommand{\arraystretch}{1.12}
	\setlength{\tabcolsep}{9pt}
	\begin{tabular}{|c|c|cc|ccc|cc|}
		\hline
		Состояние & $L_1$ & $L_2$ & $L_3$ & $L_4$ & $L_5$ & $L_6$ & $L_7$ & $L_8$ \\
		\hline
		$S_0$          & 1 & 0 & 0 & 1 & 1 & 1 & 0 & 0 \\
		$S_1$          & 1 & 0 & 0 & 0 & 0 & 1 & 0 & 0 \\
		$S_2$          & 1 & 0 & 0 & 0 & 1 & 0 & 0 & 0 \\
		$S_3$          & 1 & 0 & 0 & 1 & 0 & 0 & 0 & 0 \\
		$S_4,S_5,S_6$  & 1 & 0 & 1 & 0 & 1 & 1 & 0 & 0 \\
		$S_7$          & 1 & 1 & 0 & 1 & 1 & 0 & 1 & 0 \\
		$S_8$          & 1 & 1 & 0 & 0 & 1 & 0 & 0 & 0 \\
		$S_9$          & 1 & 1 & 0 & 1 & 0 & 0 & 0 & 0 \\
		$S_{10}$       & 0 & 0 & 0 & 1 & 1 & 1 & 0 & 0 \\
		$S_{11}$       & 1 & 1 & 1 & 1 & 0 & 0 & 0 & 0 \\
		$S_{12}$       & 1 & 0 & 0 & 1 & 1 & 1 & 0 & 1 \\
		\hline
	\end{tabular}
\end{table}

Триггеры хранят два дополнительных уровня. $L_9$ разрешает ход секундомера и изменяется командами $i_4$ и $i_5$. $L_{10}$ разрешает срабатывание будильника и переключается по $i_7$. Смена режима индикации сохраняет оба уровня.
